% =====================================================
% preamble.tex — Configuración común de estilo
% Para enunciados de ejercitación en Tekla Structures
% =====================================================

% --- Codificación e idioma ---
\usepackage[utf8]{inputenc}
\usepackage[T1]{fontenc}
\usepackage[spanish,es-tabla,es-noshorthands]{babel}
\usepackage{lmodern}
\usepackage{microtype}

% --- Geometría de página ---
\usepackage[a4paper,margin=2.5cm,headheight=14pt]{geometry}

% --- Matemática y unidades ---
\usepackage{amsmath,amssymb}
\usepackage{siunitx}
\sisetup{
  output-decimal-marker={,},
  per-mode=symbol,
  group-separator={.},
  group-minimum-digits=4
}

% --- Tablas ---
\usepackage{booktabs}
\usepackage{array}
\usepackage{multirow}
\usepackage{tabularx}
\usepackage{longtable}

% --- Figuras ---
\usepackage{graphicx}
\usepackage{caption}
\usepackage{subcaption}
\usepackage{float}

% --- Gráficos vectoriales ---
\usepackage{tikz}
\usetikzlibrary{arrows.meta,positioning,calc,patterns,shapes.geometric}

% --- Listas ---
\usepackage{enumitem}
\setlist{itemsep=2pt,parsep=2pt}

% --- Color y cajas ---
\usepackage{xcolor}
\definecolor{teklablue}{RGB}{0,87,156}
\definecolor{boxgray}{RGB}{245,245,245}
\definecolor{accentorange}{RGB}{230,120,30}

\usepackage[most]{tcolorbox}

% Caja para comandos / atajos
\newtcolorbox{comando}{
  colback=boxgray,
  colframe=teklablue,
  boxrule=0.5pt,
  arc=2pt,
  left=6pt, right=6pt, top=4pt, bottom=4pt,
  fonttitle=\bfseries,
  title=Comando Tekla
}

% Caja para tips / consejos
\newtcolorbox{tip}{
  colback=yellow!8,
  colframe=accentorange,
  boxrule=0.5pt,
  arc=2pt,
  fonttitle=\bfseries,
  title=Tip
}

% Caja para entregable
\newtcolorbox{entregable}{
  colback=green!5,
  colframe=green!50!black,
  boxrule=0.5pt,
  arc=2pt,
  fonttitle=\bfseries,
  title=Entregable
}

% Caja para advertencia
\newtcolorbox{advertencia}{
  colback=red!5,
  colframe=red!70!black,
  boxrule=0.5pt,
  arc=2pt,
  fonttitle=\bfseries,
  title=Atención
}

% --- Código y comandos en línea ---
\usepackage{listings}
\lstset{
  basicstyle=\ttfamily\small,
  breaklines=true,
  frame=single,
  backgroundcolor=\color{boxgray},
  rulecolor=\color{teklablue}
}

% Comando para resaltar comandos/menús de Tekla en línea
% Uso: \tk{File > Open}  o  \tk{Component Catalog}
\newcommand{\tk}[1]{\textsf{\textbf{\textcolor{teklablue}{#1}}}}

% Comando para nombres de archivo / rutas
\newcommand{\arch}[1]{\texttt{#1}}

% Comando para teclas / shortcuts
% Uso: \key{Ctrl+S}
\newcommand{\key}[1]{\fbox{\footnotesize\texttt{#1}}}

% --- Encabezado y pie de página ---
\usepackage{fancyhdr}
\pagestyle{fancy}
\fancyhf{}
\fancyhead[L]{\small\textit{\tituloEjercicio}}
\fancyhead[R]{\small\thepage}
\fancyfoot[C]{\small Tekla Structures 2022 — Nivel avanzado}
\renewcommand{\headrulewidth}{0.4pt}
\renewcommand{\footrulewidth}{0.4pt}

% --- Hipervínculos y referencias ---
\usepackage{hyperref}
\hypersetup{
  colorlinks=true,
  linkcolor=teklablue,
  urlcolor=teklablue,
  citecolor=teklablue,
  pdfborder={0 0 0}
}
\usepackage[noabbrev,spanish]{cleveref}

% --- Formato de secciones ---
\usepackage{titlesec}
\titleformat{\section}
  {\Large\bfseries\color{teklablue}}
  {\thesection.}{0.5em}{}
\titleformat{\subsection}
  {\large\bfseries\color{teklablue!80!black}}
  {\thesubsection.}{0.5em}{}
